
\Problem{Urgències Hospitalàries}

\Statement

Es vol simular el funcionament d'unes urgències hospitalàries.
Els pacients arriben a urgències amb una certa prioritat (1, 2 o
3, on 1 és la més alta) i un identificador. Quan el metge està
lliure, s'atén al pacient amb la prioritat més alta (el número
més petit).

\medskip

A les urgències, hi ha una sola consulta, de manera que només un
pacient pot estar dins alhora. L'entrada consisteix en una
seqüència d'events de dos tipus:

\begin{itemize}

\item \verb|arriba <id> <prioritat> <timestamp>|: arriba un
pacient amb identificador \verb|id|, prioritat \verb|prioritat| i
a l'instant \verb|timestamp|.

\item \verb|seguent <timestamp>|: el pacient que estava a la
consulta surt, i el pacient a la cua amb més prioritat entra a la
consulta a l'instant \verb|timestamp|.

\end{itemize}

\medskip

Cada cop que un pacient entra a la consulta (event
\verb|seguent|), s'escriu el seu identificador. Quan es llegeix
l'event \verb|seguent|, el pacient que entra és el que té la
prioritat més alta entre els que estan esperant. Si hi ha empat,
s'atén al que va arribar abans.

\medskip

Al final de l'entrada, s'escriu el temps d'espera mitjà dels
pacients atesos per a cada prioritat (1, 2 i 3). El temps
d'espera d'un pacient és la diferència entre el moment en què
entra a la consulta i el moment en què va arribar.

\Input

Una seqüència d'events, un per línia. Els events tenen el format
descrit més amunt. L'entrada garanteix que els events tenen
sentit: mai hi ha més d'un pacient a la consulta, i quan es fa
\verb|seguent| hi ha almenys un pacient esperant.

\Output

Per a cada event \verb|seguent|, s'escriu l'identificador del
pacient que entra a la consulta.

\medskip

Al final, per a cada prioritat que hagi tingut almenys un pacient
atès, s'escriu una línia amb el format 
\verb|prioritat <p>: <mitjana>|, on \verb|<p>| és la prioritat i 
\verb|<mitjana>| és el temps d'espera mitjà amb 2 decimals.

\Observation

Als fitxers públics (icona del gatet) trobaràs: \verb|main.cc|
(programa principal que has de completar), \verb|pacient.hh|, els
fitxers \verb|*.hh| dels contenidors de PRO2 que pots necessitar
(incloent \verb|assert.hh|), un \verb|Makefile|,
\verb|.clang-format| i el directori \verb|.vscode|.

\medskip

Has de completar el fitxer \verb|main.cc| i enviar-lo.

\medskip

Per posar la precisió de la sortida a 2 decimals, utilitza al
principi del \verb|main|:

\begin{verbatim}
cout.setf(ios::fixed);
cout.precision(2);
\end{verbatim}

\SampleTwoCol
